% English translation of capitulos/01-introducao.tex, made by Claude Code (Claude Fable 5.2) on 24 November 2026.
% Pending review by the author. When the Portuguese file changes, update this one (scripts/verificar_traducao.py).
\section{Introduction}
Large language models (LLMs) have become part of software development tools. \citet{chen2021codex}, for
example, reported in 2021 the use of a production version of Codex in GitHub Copilot. As a result,
part of a programmer's work becomes writing instructions for the model, defining
context, examples and constraints. This practice is known as prompt engineering.

The literature on the subject is growing fast and uses techniques, models, tasks and metrics that differ
widely from one another. In the search carried out for this work, the same structure of terms retrieved 2,304
unique studies in three databases (Section~\ref{sec:busca}). This volume makes it hard to know what has already
been evaluated, under which conditions and with which limits. The motivation of this work is the expectation, which it
does not test, that an organized synthesis helps developers and researchers choose techniques
more judiciously.

The scope is restricted to pre-trained models used through prompts in four tasks: code generation,
translation, refactoring and debugging, including program repair. Training, fine-tuning
and model architecture are not covered, except when the study also evaluates prompting.
This work does not propose models and does not run experiments of its own to measure the effectiveness of
prompting techniques in the four tasks. LLMs are used and evaluated only in the operation of the review
(Section~\ref{sec:metodologia}), and those numbers do not answer the research questions.

\subsection{Objectives}
The objective is to organize evidence on prompt engineering techniques applied to code generation,
translation, refactoring and debugging, by means of a literature review and a
taxonomy to support developers. The specific objectives are:
\begin{itemize}
  \item to identify which prompting techniques are evaluated and what evidence of effectiveness is
  presented, under which models, tasks, metrics and conditions;
  \item to identify the technical risks reported and the mitigation strategies proposed or evaluated;
  \item to relate techniques, tasks and conditions in a taxonomy and in a decision-support table,
  examined in eight scenarios for clarity, grounding in the sources and concrete guidance.
\end{itemize}
This examination does not measure the effectiveness of the techniques or the productivity of those who use them.

\subsection{Organization of the paper}
Section~\ref{sec:fundamentacao} presents the concepts and the related work.
Section~\ref{sec:metodologia} describes the protocol, what was carried out and what is left for the final version. The following sections bring together the results, the discussion, with the
application that makes the decision-support table available for consultation (Section~\ref{sec:aplicacao}), and the conclusion.
